\section{Concluding Discussion}
\label{sec:res-disc}

\xiaoyan{Discuss more on the linear probing results and the mismatch between causality test and the probe.}
\anita{Implications paragraph add later}

Taken together, the attention and probing results suggest that chemical information is not represented uniformly across mechanisms or layers in these models. Attention exhibits only weak and highly selective alignment with molecular geometry, indicating that spatial structure is not encoded through a small set of geometry-tracking heads. In contrast, probing and causal intervention reveal that chemically meaningful properties are strongly represented within the residual stream, but become computationally usable at substantially different depths across models. More broadly, our results move beyond asking whether chemical information is present, toward understanding how chemically relevant structure is internally organized and used during computation.

\para{The primary difference between \molm and \smi lies in when chemical information becomes functionally usable.}
Our results suggest that \smi makes contextual chemical information functionally usable substantially earlier in the encoder, whereas \molm often does so only near the top layers. This distinction may matter most for downstream tasks requiring multi-step integration of chemical information. Properties relevant to ADMET prediction, reactivity, and molecular binding often depend on combining multiple local and non-local chemical cues rather than detecting a single motif. When a representation becomes usable early, subsequent layers can build on it as an intermediate computational feature for forming more global or task-specific representations. By contrast, representations that become usable only near the top of the encoder may still support final prediction, but have less opportunity to influence preceding computation.

This distinction may be particularly important in practical transfer settings, where molecular encoders are commonly used through fixed embeddings, shallow predictors, parameter-efficient adaptation, partial freezing, or intermediate-layer readouts rather than fully end-to-end optimization~\cite{honda2019smiles,rollins2024molprop,cai2026chemfm,harod2025effichem,pinto2025unlocking}. Under these constraints, downstream performance may depend less on what information is eventually present in the model and more on which chemical features are already accessible, and at what depth. Layerwise functional accessibility may therefore be an important predictor of downstream robustness, transferability, and data efficiency.

% \para{These findings also suggest design principles for better chemistry foundation models.}
% First, tokenizer-aware evaluation is essential. Some properties appear at layer~0 because they are already specified by token identity, and these cases should not be treated as evidence of learned chemical abstraction. More informative benchmarks are contextual properties that require integration beyond the local token. Second, pretraining should be evaluated not only by final task accuracy, but also by whether it produces simple and persistent internal features that remain usable across depth. This perspective favors objectives and architectures that encourage early integration of relational structure rather than deferring chemically meaningful computation to the final layers. Third, the weak attention-distance correlations highlight a limitation of standard SMILES-only pretraining: strong predictive performance does not necessarily imply a spatially grounded internal representation. If spatial reasoning is the target capability, stronger geometric inductive bias or multimodal 3D supervision may be required.

\para{Future Directions.}
One natural next step is to study the internal geometry of chemical representations more directly. While linear probing identifies which properties are recoverable, it does not characterize how different chemical features are arranged relative to one another within representation space, whether they occupy shared or independent subspaces, or how these structures evolve across layers. Future work could examine feature overlap, representation superposition, and layerwise organization using methods such as sparse autoencoders, subspace analysis, or representation similarity techniques.

A second direction is to trace how chemical information flows through the network during computation. Our causal interventions show when a representation becomes functionally usable, but do not yet identify which downstream components consume that information or how multiple chemical features are composed into higher-level reasoning. Studying attention circuits, residual-stream interactions, and feature propagation across layers may help reveal how models combine local chemical descriptors into more global molecular representations.

Finally, an important open question is how different pretraining choices shape the resulting internal representations. Although \molm{} and \smi{} achieve similar downstream performance, they differ substantially in when and how chemical information becomes usable. Extending this analysis to additional architectures, objectives, and training datasets may help identify which design choices lead to more structured, transferable, or computationally efficient chemical representations.
